\documentclass[11pt,a4paper]{article}
\usepackage[margin=25mm,headheight=14pt]{geometry}
\usepackage{fontspec}
\setmainfont{lmroman10-regular.otf}[BoldFont=lmroman10-bold.otf,ItalicFont=lmroman10-italic.otf,BoldItalicFont=lmroman10-bolditalic.otf]
\setsansfont{lmsans10-regular.otf}
\setmonofont{lmmono10-regular.otf}[Scale=MatchLowercase]
\newfontfamily\authorsymbol{Noto Sans Symbols}[Scale=MatchUppercase]
\usepackage{amsmath,amssymb,amsthm,mathtools}
\usepackage{booktabs,array,longtable}
\usepackage{microtype,needspace}
\usepackage[hidelinks,unicode]{hyperref}
\usepackage{fancyhdr}
\pagestyle{fancy}\fancyhf{}
\fancyhead[L]{\small\itshape Twenty signed products for a composition tensor}
\fancyhead[R]{\small\thepage}
\renewcommand{\headrulewidth}{0pt}
\setlength{\parindent}{1.2em}\setlength{\parskip}{0pt}
\setlength{\emergencystretch}{2em}
\newtheorem{theorem}{Theorem}
\newtheorem{proposition}[theorem]{Proposition}
\newtheorem{lemma}[theorem]{Lemma}
\newtheorem{corollary}[theorem]{Corollary}
\DeclareMathOperator{\rank}{rank}
\DeclareMathOperator{\rad}{rad}
\DeclareMathOperator{\supp}{supp}
\newcommand{\R}{\mathbb R}\newcommand{\Q}{\mathbb Q}
\newcommand{\Z}{\mathbb Z}\newcommand{\ind}{\mathbf1}
\newcommand{\rkplus}{\operatorname{rank}_{+}}
\newcommand{\rkreal}{\operatorname{rank}_{\R}}
\hypersetup{pdftitle={Twenty Signed Products for a Twelve-State Composition Tensor},pdfauthor={lmtlssss🜐},pdfsubject={Mathematics: tensor rank and algebraic complexity},pdfkeywords={tensor rank, nonnegative rank, finite monoids, bilinear complexity}}
\title{\LARGE Twenty signed products for a\protect\\twelve-state composition tensor}
\author{lmtlssss{\authorsymbol 🜐}}
\date{27 October 2026}
\begin{document}
\maketitle\thispagestyle{plain}
\begin{abstract}
We give a twenty-term signed decomposition of the multiplication tensor of a
specified twelve-element transformation monoid. Its nonnegative rank is
exactly 35. An integral change of coordinates separates five scalar
characters, a two-by-two matrix product, and three additional bilinear
forms. Strassen's seven products for the matrix component give twenty
products in total. A square-zero radical and the semisimple quotient
$\R^5\times M_2(\R)$ give the lower bound 18. For every positive integer
$q$, the nonnegative rank of the tensor power is $35^q$, whereas its
signed rank is at most $20^q$. The coordinate formulas and exact
integer certificates are supplied.
\end{abstract}
\noindent\textit{Mathematics Subject Classification (2020).}
15A69; 16Z05, 20M20.\par
\noindent\textit{Keywords.} Tensor rank; finite monoid; bilinear complexity;
nonnegative decomposition.\par
\noindent\textit{Version.} 1.0.0.\quad
\textit{DOI.} \href{https://doi.org/10.5281/zenodo.23695066}{10.5281/zenodo.23695066}.

\section{The result}

Let $C$ be the composition tensor of the monoid $S$ defined below.
The preceding note~\cite{positive} determined its nonnegative rank:
$\rkplus(C)=35$. Here we allow negative coefficients and construct the
same bilinear map with twenty products. The change is in the
representation of the tensor, with its coefficients fixed.

For a real three-way tensor, $\rkreal(C)$ is the least $r$ such that
\begin{equation}\label{eq:rank}
 C=\sum_{\ell=1}^{r}\alpha_\ell\otimes\beta_\ell\otimes\gamma_\ell,
 \qquad \alpha_\ell,\beta_\ell,\gamma_\ell\in\R^{12}.
\end{equation}
The nonnegative rank requires all three factors to be entrywise
nonnegative. A decomposition~\eqref{eq:rank} computes the associated
bilinear map with $r$ products of linear forms. Linear operations belong
to a separate cost account.

\begin{theorem}\label{thm:main}
The tensor $C$ has an integer twenty-term decomposition, with every
factor coefficient in $\{0,\pm1,\pm2\}$. Moreover,
\begin{equation}\label{eq:bounds}
 18\leq\rkreal(C)\leq20,\qquad \rkplus(C)=35.
\end{equation}
For every integer $q\geq1$,
\begin{equation}\label{eq:power}
 \rkplus(C^{\otimes q})=35^q,
 \qquad \rkreal(C^{\otimes q})\leq20^q.
\end{equation}
\end{theorem}

The construction uses Strassen's matrix algorithm~\cite{strassen}.
The lower bound uses the Alder--Strassen inequality~\cite{alder}, as
formulated for bilinear rank in Bl\"aser~\cite{blaser}.
This note records the explicit construction for $S$, its algebra
structure, and the resulting tensor-power separation.

\section{The finite monoid}

A map $f:\{0,1,2,3\}\to\{0,1,2,3\}$ is written as the tuple
$(f(0),f(1),f(2),f(3))$. Consider the following twelve maps, in the given
order:
\begin{equation}\label{eq:maps}
\begin{array}{c|cccccc}
 i&0&1&2&3&4&5\\\hline
 m_i&(0,0,0,0)&(0,0,0,1)&(0,1,1,1)&(0,1,2,2)&(0,1,2,3)&(1,1,1,1)\\[3pt]
 i&6&7&8&9&10&11\\\hline
 m_i&(1,1,2,2)&(1,1,2,3)&(2,2,2,2)&(2,2,2,3)&(2,3,3,3)&(3,3,3,3).
\end{array}
\end{equation}
Set $m_i*m_j=m_j\circ m_i$ and write $M(i,j)$ for its index.
Table~\ref{tab:monoid} proves closure. Associativity follows from function
composition, and $m_4$ is the identity. The six maps with indices
$1,2,3,7,9,10$ generate $S$.

\begin{table}[ht]
\centering\small\setlength{\tabcolsep}{5.2pt}
\begin{tabular}{r|rrrrrrrrrrrr}
\toprule
$M(i,j)$&0&1&2&3&4&5&6&7&8&9&10&11\\\midrule
0 & 0 & 0 & 0 & 0 & 0 & 5 & 5 & 5 & 8 & 8 & 8 & 11 \\
1 & 0 & 0 & 1 & 1 & 1 & 5 & 5 & 5 & 8 & 8 & 9 & 11 \\
2 & 0 & 0 & 2 & 2 & 2 & 5 & 5 & 5 & 8 & 8 & 10 & 11 \\
3 & 0 & 0 & 2 & 3 & 3 & 5 & 6 & 6 & 8 & 8 & 10 & 11 \\
4 & 0 & 1 & 2 & 3 & 4 & 5 & 6 & 7 & 8 & 9 & 10 & 11 \\
5 & 0 & 0 & 5 & 5 & 5 & 5 & 5 & 5 & 8 & 8 & 11 & 11 \\
6 & 0 & 0 & 5 & 6 & 6 & 5 & 6 & 6 & 8 & 8 & 11 & 11 \\
7 & 0 & 1 & 5 & 6 & 7 & 5 & 6 & 7 & 8 & 9 & 11 & 11 \\
8 & 0 & 0 & 5 & 8 & 8 & 5 & 8 & 8 & 8 & 8 & 11 & 11 \\
9 & 0 & 1 & 5 & 8 & 9 & 5 & 8 & 9 & 8 & 9 & 11 & 11 \\
10 & 0 & 2 & 5 & 8 & 10 & 5 & 8 & 10 & 8 & 10 & 11 & 11 \\
11 & 0 & 5 & 5 & 8 & 11 & 5 & 8 & 11 & 8 & 11 & 11 & 11 \\
\bottomrule
\end{tabular}
\caption{The product applies the row map first, then the column map.}
\label{tab:monoid}
\end{table}

Let $A=\R[S]$, with basis $e_0,\ldots,e_{11}$ and product
$e_ie_j=e_{M(i,j)}$. For coefficient columns $a,b\in\R^{12}$, denote
by $p=a*b$ the coefficient column of their product. Thus
\begin{equation}\label{eq:product}
 C_{ijk}=\ind_{\{M(i,j)=k\}},\qquad
 p_k=\sum_{M(i,j)=k}a_ib_j.
\end{equation}
There are exactly 144 nonzero tensor entries. The indexing and product
order agree with the preceding nonnegative-rank calculation~\cite{positive}.

\section{Coordinates for the twenty products}\label{sec:coordinates}

Define the linear coordinate change $u=Qa$ by
\begin{equation}\label{eq:Q}
\begin{alignedat}{2}
 u_0&=a_4, &\qquad u_6&=a_1,\\
 u_1&=a_4+a_7, &u_7&=a_4+a_7+a_9,\\
 u_2&=a_3+a_4, &u_8&=a_2+a_3+a_4,\\
 u_3&=a_3+a_4+a_6+a_7, &u_9&=a_8+a_9+a_{10}+a_{11},\\
 u_4&=\sum_{i=0}^{11}a_i, &u_{10}&=a_{11},\\
 u_5&=a_{10}, &u_{11}&=a_0.
\end{alignedat}
\end{equation}
For any column $z=Qp$, the inverse is
\begin{equation}\label{eq:Qinv}
\begin{alignedat}{2}
 p_0&=z_{11}, &\qquad p_6&=z_0-z_1-z_2+z_3,\\
 p_1&=z_6, &p_7&=z_1-z_0,\\
 p_2&=z_8-z_2, &p_8&=z_1-z_5-z_7+z_9-z_{10},\\
 p_3&=z_2-z_0, &p_9&=z_7-z_1,\\
 p_4&=z_0, &p_{10}&=z_5,\\
 p_5&=z_2-z_3+z_4-z_6-z_8-z_9-z_{11},
 &p_{11}&=z_{10}.
\end{alignedat}
\end{equation}
Substitution proves that these maps are mutual inverses over $\Z$.
In particular, $Q$ is unimodular.

For $v=Qb$ and $z=Q(a*b)$, the first five coordinates satisfy
\begin{equation}\label{eq:characters}
 z_i=u_iv_i\qquad(0\leq i\leq4).
\end{equation}
The next four coordinates form the matrix product
\begin{equation}\label{eq:matrix}
 \begin{pmatrix}z_8&z_6\\z_5&z_7\end{pmatrix}
 =\begin{pmatrix}v_8&v_6\\v_5&v_7\end{pmatrix}
  \begin{pmatrix}u_8&u_6\\u_5&u_7\end{pmatrix}.
\end{equation}
The remaining coordinates are
\begin{align}
 z_9&=u_9v_3+u_4v_9,\label{eq:z9}\\
 z_{10}&=u_{10}v_7+(u_4-u_6-u_8-u_{11})v_5+u_4v_{10},\label{eq:z10}\\
 z_{11}&=u_4v_{11}+(u_4-u_5-u_7-u_{10})v_6+u_{11}v_8.\label{eq:z11}
\end{align}

\begin{lemma}\label{lem:coordinates}
Equations~\eqref{eq:characters}--\eqref{eq:z11} hold for every $a,b$.
\end{lemma}
\begin{proof}
Both sides are bilinear. For $a=e_i$ and $b=e_j$, the left side is
$Qe_{M(i,j)}$. Substitution of the corresponding columns of $Q$ into
the right side gives the same column for each entry of
Table~\ref{tab:monoid}. This finite identity proves the assertion by
bilinearity. The supplied verifier checks all 144 basis pairs with
integer arithmetic. It also checks the expanded twenty-term tensor
identity independently of these coordinate formulas.
\end{proof}

To evaluate~\eqref{eq:matrix}, use seven scalar products:
\begin{equation}\label{eq:strassen}
\begin{alignedat}{2}
 s_1&=(u_8+u_7)(v_8+v_7),&\qquad s_2&=u_8(v_5+v_7),\\
 s_3&=(u_6-u_7)v_8,&s_4&=(u_5-u_8)v_7,\\
 s_5&=u_7(v_8+v_6),&s_6&=(u_8+u_6)(v_5-v_8),\\
 s_7&=(u_5+u_7)(v_6-v_7).&&
\end{alignedat}
\end{equation}
Then
\begin{equation}\label{eq:strassenout}
 z_8=s_1+s_4-s_5+s_7,\quad z_6=s_3+s_5,\quad
 z_5=s_2+s_4,\quad z_7=s_1-s_2+s_3+s_6.
\end{equation}
These are Strassen's formulas~\cite{strassen} with the order of the
matrix factors fixed by~\eqref{eq:matrix}.

\begin{proposition}\label{prop:upper}
Equations~\eqref{eq:Q}--\eqref{eq:strassenout} give an exact algorithm
with twenty products of an $a$-linear form and a $b$-linear form.
\end{proposition}
\begin{proof}
Equation~\eqref{eq:characters} uses five products.
Equations~\eqref{eq:z9}--\eqref{eq:z11} use $2+3+3=8$ products.
Equation~\eqref{eq:strassen} uses seven. The coordinate changes and
reconstruction use only linear operations. Therefore the total is
$5+8+7=20$. Expansion in the original coordinates gives
\begin{equation}\label{eq:certificate}
 C_{ijk}=\sum_{\ell=1}^{20}
 \alpha_{\ell i}\beta_{\ell j}\gamma_{\ell k}
 \quad(0\leq i,j,k<12).
\end{equation}
The coefficients in the expansion belong to $\{0,\pm1,\pm2\}$.
The supplied certificate lists all 720 factor coefficients.
The identity holds over $\Z$, hence over $\Q$, $\R$ and $\mathbb C$.
\end{proof}

\section{The lower bound}

The first five coordinate functions in~\eqref{eq:Q} are unital algebra
homomorphisms from $A$ to $\R$, by~\eqref{eq:characters}. Write them as
$\chi_0,\ldots,\chi_4$. Define
\begin{equation}\label{eq:rho}
 \rho(a)=\begin{pmatrix}u_8&u_6\\u_5&u_7\end{pmatrix}.
\end{equation}
Equation~\eqref{eq:matrix} states that $\rho(a*b)=\rho(b)\rho(a)$.
Thus $a\mapsto\rho(a)^{\mathsf T}$ is a unital homomorphism.

\begin{lemma}\label{lem:radical}
The algebra $A$ has a three-dimensional radical $J$ with $J^2=0$, and
\begin{equation}\label{eq:quotient}
 A/J\cong\R^5\times M_2(\R).
\end{equation}
\end{lemma}
\begin{proof}
Consider the homomorphism
\[
 \Phi(a)=\bigl(\chi_0(a),\ldots,\chi_4(a),\rho(a)^{\mathsf T}\bigr).
\]
Its nine coordinate rows are precisely the first nine rows of $Q$,
up to order. They are independent because $Q$ is invertible, so
$\Phi$ is surjective. Its kernel is
\begin{equation}\label{eq:J}
 J=\operatorname{span}\{e_5-e_0,\ e_8-e_0,\ e_{11}-e_0\}.
\end{equation}
These three independent vectors vanish under $\Phi$, and the kernel
has dimension three.

The maps $m_0,m_5,m_8,m_{11}$ are constant. Under our product convention,
$e_re_s=e_s$ for their indices. Hence
$(e_r-e_0)(e_s-e_0)=e_s-e_0-e_s+e_0=0$.
This proves $J^2=0$. Every nilpotent ideal belongs to the Jacobson
radical. Conversely, $A/J$ is the semisimple algebra in
\eqref{eq:quotient}, so $\rad A\subseteq J$. Therefore $J=\rad A$.
\end{proof}

\begin{proposition}\label{prop:lower}
The real bilinear rank of $C$ is at least 18.
\end{proposition}
\begin{proof}
The quotient in~\eqref{eq:quotient} has six simple factors, so $A$ has
six maximal two-sided ideals. The Alder--Strassen bound gives
\[
 R(A)\geq2\dim A-t(A)=2\cdot12-6=18.
\]
Here $R(A)$ is the rank of the multiplication tensor of $A$, and
$t(A)$ is its number of maximal two-sided ideals~\cite{alder,blaser}.
Changing the algebra basis preserves this rank. In the basis $e_i$,
its multiplication tensor is $C$.
\end{proof}

Together, Propositions~\ref{prop:upper} and~\ref{prop:lower} give
$\rkreal(C)\in\{18,19,20\}$. The construction specifies an upper bound
of twenty. Determining which of these three values is attained remains
an open calculation for this algebra.

\section{Tensor powers and the nonnegative bound}

The proof of the positive-rank theorem in~\cite{positive} gives two
finite certificates. Thirty-five monochromatic rectangles partition
all 144 nonzero entries of $C$. A set $W$ of 35 nonzero triples is
pairwise incompatible: for any two distinct members, at least one
coordinate-wise mixed triple is a zero of $C$. Both certificates are
included in the present release.

For completeness, a nonnegative rank-one term has support of the form
$I\times J\times K$. In a nonnegative decomposition of $C$, this box
must lie inside $\supp C$. A box therefore covers at most one member
of $W$. This gives the lower bound 35, and the rectangle partition
gives the matching upper bound.

\Needspace{8\baselineskip}
\begin{proposition}\label{prop:powers}
For every $q\geq1$, $\rkplus(C^{\otimes q})=35^q$.
\end{proposition}
\begin{proof}
Tensor the 35-term nonnegative decomposition to obtain $35^q$ terms.
For the reverse inequality, group each tensor leg into a $q$-tuple.
The Cartesian product $W^q$ has $35^q$ nonzero entries. Two distinct
members differ in at least one coordinate. In that coordinate,
incompatibility of the corresponding members of $W$ supplies a mixed
zero. Apply the same choices of tensor legs to the two $q$-tuples.
The resulting entry of $C^{\otimes q}$ has a zero factor, so it is zero.
Thus $W^q$ is pairwise incompatible. Every nonnegative rank-one term
covers at most one of its entries. At least $35^q$ terms are required.
\end{proof}

Tensoring~\eqref{eq:certificate} gives a signed decomposition with
$20^q$ terms. Therefore
\begin{equation}\label{eq:separation}
 \frac{\rkplus(C^{\otimes q})}{\rkreal(C^{\otimes q})}
 \geq\left(\frac{35}{20}\right)^q
 =\left(\frac74\right)^q.
\end{equation}
This proves Theorem~\ref{thm:main}.

The operation $C^{\otimes q}$ has two input vectors of length $12^q$
and one output vector of length $12^q$. It is a tensor product of
bilinear maps. Repeated multiplication of $q$ elements in the fixed
algebra $A$ is a different operation. Equation~\eqref{eq:separation}
compares bilinear product counts for the tensor product.

\section{Certificates and reproduction}

The release contains \texttt{certificate.json} and a standard-library
Python verifier. The certificate includes the indexed maps, table,
coordinate matrix, twenty signed terms, positive rectangle partition,
and lower-bound witnesses. The verifier reconstructs composition from
the maps and checks all 1,728 coefficient equalities in
\eqref{eq:certificate} over the integers.

A separately written implementation of
\eqref{eq:Q}--\eqref{eq:strassenout} agrees with the table on all
144 input basis pairs. These comparisons determine the complete
bilinear map. Rational elimination checks the nine-dimensional quotient
and the kernel in~\eqref{eq:J}. The verifier also checks all 595 base
witness pairs and all 749,700 pairs in the tensor-square witness set.
The general tensor-power result follows from
Proposition~\ref{prop:powers}, rather than a finite extrapolation.

The following command reproduces the certificate checks:
\begin{center}\texttt{python3 verify.py}\end{center}
The written proofs establish the mathematical statements. The executable
checks provide exact, inspectable evidence for the finite tables and
formulas. The algebraic rank model counts products of linear forms.
Addition counts, storage, numerical stability and execution time require
their own analysis.

\begin{thebibliography}{9}
\bibitem{positive}
lmtlssss{\authorsymbol 🜐}.
\emph{Exact nonnegative rank of a twelve-state composition tensor}.
Version 1.0.0, 2026.
\href{https://doi.org/10.5281/zenodo.23678773}{doi:10.5281/zenodo.23678773}.

\bibitem{strassen}
V.~Strassen.
Gaussian elimination is not optimal.
\emph{Numerische Mathematik} \textbf{13} (1969), 354--356.
\href{https://doi.org/10.1007/BF02165411}{doi:10.1007/BF02165411}.

\bibitem{alder}
A.~Alder and V.~Strassen.
On the algorithmic complexity of associative algebras.
\emph{Theoretical Computer Science} \textbf{15} (1981), 201--211.
\href{https://doi.org/10.1016/0304-3975(81)90070-0}{doi:10.1016/0304-3975(81)90070-0}.

\bibitem{blaser}
M.~Bl\"aser.
A complete characterization of the algebras of minimal bilinear complexity.
\emph{SIAM Journal on Computing} \textbf{34} (2005), 277--298.
\href{https://doi.org/10.1137/S0097539703438277}{doi:10.1137/S0097539703438277}.
\end{thebibliography}
\end{document}
